\documentclass[11pt,letterpaper]{article}
\usepackage{physical_ai_legislation}
\usepackage[backend=biber,style=numeric,sorting=none]{biblatex}
\addbibresource{physical_ai_patient_rights.bib}

\title{%
  \textbf{Assembly Bill 2847}\\[6pt]
  \large California Physical AI Patient Rights\\
  and Robotic Safety Act of 2026\\[6pt]
  \normalsize An act to add Article 7.8 (commencing with Section 1368.10)\\
  to Chapter 2.2 of Division 2 of the Health and Safety Code,\\
  relating to patient rights in physical AI oncology clinical trials.
}
\author{%
  Kevin Kawchak\\
  CEO, ChemicalQDevice\\
  \texttt{kevink@chemicalqdevice.com}
}
\date{March 24, 2026}

\begin{document}
\maketitle
\thispagestyle{fancy}

\begin{center}
\footnotesize
\textit{Unofficial independent draft derived from prior legislation;
no third-party endorsement, sponsorship, affiliation, or authorization
is expressed or implied; and was adapted using Claude Code Opus 4.6.}
\end{center}

\vspace{0.5em}

\begin{center}
\footnotesize
\textit{This draft is independent and is not endorsed, sponsored, or approved
by any trial sponsor, CRO, site, IRB, regulator, or medical society;
and was adapted using Claude Code Opus 4.6.}
\end{center}

\vspace{1em}
\tableofcontents
\newpage

\section{Legislative Findings and Declarations}

The Legislature finds and declares all of the following:

\begin{enumerate}[label=(\alph*)]
\item Physical AI oncology clinical trials represent a fundamental shift in the
  relationship between patients and clinical care systems, in which control has
  now shifted towards the patient's side, with more patient rights than
  traditional clinical trial formats provide.

\item The 24-hour on-demand Physical AI oncology trial simulation demonstrated
  that patients can choose their own scheduling across a full 24-hour period,
  experience average wait times of 8 minutes, and achieve patient satisfaction
  scores of 4.7 out of 5.0, compared to the rigid scheduling and extended
  wait times typical of traditional oncology trials~\cite{newtrial2026}.

\item More patients will be treated with more robots and fewer workers, with the
  simulation demonstrating 168 patients served by 29 robot instances staffed
  by 8 to 10 human full-time equivalents across three shifts, compared to 5 to
  15 patients per day in traditional formats requiring 80 to 120
  staff~\cite{formatcomparison2026}.

\item The cancer patient journey through a regulated Physical AI oncology trial
  documented a complete patient trajectory from pre-screening through data
  lock across 10 stages, achieving complete Response, R0 resection, and
  reduction in recurrence risk from 35 percent to 3 percent, demonstrating
  that patient outcomes improve under Physical AI
  care~\cite{patientjourney2026}.

\item Pediatric patients require special protections, including age-appropriate
  demonstrations of robotic systems, parental consent provisions, and dedicated
  companion robot interactions, as established in the adapted 21 CFR Part 50
  \S~50.50 through \S~50.56~\cite{cfr50adapt}.

\item Patient-facing robot instructions for all 10 robot types have been
  developed and validated, providing step-by-step guidance for patients
  interacting with surgical robots, cobots, radiotherapy positioning systems,
  needle-placement systems, social companions, humanoids, motion-tracking
  systems, imaging assistants, steerable needles, and rehabilitation
  exoskeletons~\cite{patientinstructions2026}.

\item The simulation demonstrated that both specific details by the minute and
  details over 24 hours were handled correctly by AI, with minute-level vital
  sign monitoring, real-time adverse event detection and resolution, and
  sustained clinical operations across all 24 hours.

\item Seven adverse events occurred across 168 patients (4.2 percent incidence
  rate), all Grade 1 or Grade 2, all resolved within the same hour,
  demonstrating a safety profile that supports expanding patient access to
  Physical AI oncology care~\cite{newtrial2026}.
\end{enumerate}

\section{Patient Bill of Rights for Physical AI\\Oncology Clinical Trials}

\subsection{Right to On-Demand Scheduling}

\begin{enumerate}[label=(\alph*)]
\item Every patient enrolled in a Physical AI oncology clinical trial at an
  authorized California site shall have the right to schedule procedures,
  consultations, and follow-up visits at any time during the site's 24-hour
  operating period, 7 days per week, subject only to clinical safety
  constraints and equipment availability.

\item The site shall provide scheduling systems that enable patients to select
  appointment times without requiring approval from clinical staff, with the
  Physical AI systems managing patient flow and resource allocation
  autonomously.

\item Maximum wait times from arrival to procedure initiation shall not exceed
  30 minutes under normal operating conditions. The simulation demonstrated
  an average wait time of 8 minutes across 168
  patients~\cite{newtrial2026}.
\end{enumerate}

\subsection{Right to Full Disclosure}

\begin{enumerate}[label=(\alph*)]
\item Before any Physical AI-assisted procedure, the patient shall receive clear
  and complete disclosure of the following:
  \begin{enumerate}[label=(\arabic*)]
  \item The specific Physical AI system or systems that will be involved in
    their care, identified by type, model, and the role each system will
    perform (primary, secondary, or ancillary).
  \item Whether the system will operate in autonomous, assistive, or
    shared-autonomy mode.
  \item The system's current USL score presented in plain language, with an
    explanation of what the score means for the patient's
    procedure~\cite{usl2026}.
  \item The system's safety record, including any prior adverse events involving
    the specific system instance.
  \item Emergency stop provisions and how they will be activated if needed.
  \item The identity and qualifications of the human oversight personnel
    responsible for supervising the procedure.
  \item Simulation validation results relevant to the patient's specific
    procedure type.
  \item Digital twin usage and the patient's right to request deletion of their
    digital twin data at any time.
  \end{enumerate}

\item Disclosure materials shall be available in the patient's preferred
  language and in formats accessible to patients with visual, auditory, or
  cognitive impairments.

\item The site shall provide patient-facing robot instructions consistent with
  those developed for each of the 10 Physical AI system
  categories~\cite{patientinstructions2026}.
\end{enumerate}

\subsection{Right to Human Intervention}

\begin{enumerate}[label=(\alph*)]
\item A patient shall have the right to request human intervention or oversight
  at any point during a Physical AI-assisted procedure, and the site shall be
  staffed to accommodate such requests within 2 minutes of the request.

\item A patient may withdraw consent for Physical AI involvement at any time and
  request that the remainder of their procedure be completed by human
  clinicians, to the extent that such transition is clinically safe.

\item The site shall maintain at least one on-call emergency physician and one
  on-call pharmacist at all times during the 24-hour operating period, in
  addition to the site safety officers present on each shift.

\item Emergency stop mechanisms shall be accessible to both clinical staff and
  patients where clinically appropriate, with clear instructions provided in
  advance.
\end{enumerate}

\subsection{Right to Privacy and Data Control}

\begin{enumerate}[label=(\alph*)]
\item All patient data collected, processed, or stored by Physical AI systems
  shall be protected in accordance with HIPAA, the California Consumer Privacy
  Act, the California Privacy Rights Act, and the data governance requirements
  of the adapted ICH E6(R3)~\cite{iche6r3adapt}.

\item Patient data processed through MCP-PAI servers shall use deny-by-default
  authorization policies, meaning no data access is permitted unless
  explicitly authorized for a specific purpose~\cite{iche6r3adapt}.

\item Patients shall have the right to request a complete record of all
  Physical AI system interactions involving their care, including robotic
  telemetry logs, AI decision logs, and digital twin state data.

\item Patients shall have the right to request deletion of their digital twin
  and associated patient-specific data from Physical AI system local storage,
  subject to regulatory data retention requirements under 21 CFR Part 312.

\item Hash-chained audit trails (SHA-256) shall be maintained for all patient
  data transactions, providing an immutable record of data access and
  modification~\cite{iche6r3adapt}.
\end{enumerate}

\subsection{Right to Equitable Access}

\begin{enumerate}[label=(\alph*)]
\item Physical AI oncology clinical trial sites shall not discriminate in
  enrollment, scheduling, or treatment on the basis of race, ethnicity, sex,
  gender identity, sexual orientation, age, disability, socioeconomic status,
  insurance status, or geographic location.

\item The on-demand 24/7 scheduling format shall accommodate patients who work
  non-traditional hours, have caregiving responsibilities, or face
  transportation barriers, consistent with the patient-centric design
  demonstrated in the simulation~\cite{formatcomparison2026}.

\item The site shall provide transportation assistance, including accessible
  parking, patient drop-off zones, and coordination with public transit and
  ride-share services.

\item The per-patient cost reduction of 75 to 85 percent demonstrated in the
  simulation format comparison shall be reflected in reduced financial barriers
  for patients~\cite{formatcomparison2026}.
\end{enumerate}

\subsection{Pediatric Patient Protections}

\begin{enumerate}[label=(\alph*)]
\item Pediatric patients (under 18 years of age) shall receive age-appropriate
  explanations and demonstrations of all Physical AI systems involved in their
  care, consistent with the assent requirements of the adapted 21 CFR Part 50
  \S~50.50 through \S~50.56~\cite{cfr50adapt}.

\item Social companion robots used with pediatric patients shall follow the
  interaction protocols established in the patient robot instructions,
  including gentle high-fives, interactive games, calm-breathing prompts, and
  virtual sticker rewards~\cite{patientinstructions2026}.

\item Humanoid robots conducting rehabilitation exercises with pediatric
  patients shall follow established protocols including slow count-in and
  count-out breathing, posture cues, and interaction durations appropriate for
  the child's age and condition~\cite{patientinstructions2026}.

\item A parent or guardian shall be present or immediately available during all
  Physical AI-assisted procedures involving pediatric patients, and an advocate
  shall be appointed for wards of the state consistent with \S~50.56 of the
  adapted regulations~\cite{cfr50adapt}.

\item The simulation demonstrated successful management of pediatric patients
  including pediatric ALL (acute lymphoblastic leukemia) and AML (acute
  myeloid leukemia) cases, with companion robot-assisted anxiety reduction
  achieving 95 percent effectiveness~\cite{newtrial2026}.
\end{enumerate}

\section{Robotic Safety Standards}

\subsection{Physical AI System Safety Certification}

\begin{enumerate}[label=(\alph*)]
\item Before deployment at an authorized California site, each Physical AI
  system shall obtain safety certification through a process administered by
  the department, demonstrating the following:
  \begin{enumerate}[label=(\arabic*)]
  \item Completion of Phase 0 simulation validation across two or more
    validated simulation frameworks with cross-framework trajectory deviation
    below 2 millimeters and force discrepancy below 0.5
    Newtons~\cite{iche6r3adapt}.
  \item Achievement of minimum USL scores: 7.0 for surgical robots, 6.0 for
    therapeutic positioning and diagnostic systems, 4.0 for rehabilitative
    systems, and 3.0 for companion systems~\cite{cfr50adapt}.
  \item Completion of failure mode analysis across a minimum of 1,000
    simulated procedures~\cite{cfr312adapt}.
  \item Functional emergency stop mechanisms tested and verified.
  \item Force-limiting capabilities appropriate to the system type, including
    maximum 250 millimeters per second approach speed and force-limited contact
    for cobots~\cite{patientinstructions2026}.
  \end{enumerate}

\item Safety certification shall be renewed annually and following any
  significant hardware or software modification.
\end{enumerate}

\subsection{Continuous Safety Monitoring}

\begin{enumerate}[label=(\alph*)]
\item Each Physical AI system shall maintain continuous telemetry recording
  during all patient interactions, as demonstrated in the hourly robot logs of
  the 24-hour simulation where all 29 robot instances were tracked with
  active and standby status, performance metrics, and state
  transitions~\cite{newtrial2026}.

\item Real-time safety monitoring shall include automated detection of force
  limit exceedances, positioning errors beyond specified tolerances, AI model
  performance degradation, and cybersecurity anomalies.

\item The simulation demonstrated 99.7 percent fleet uptime across 29 robot
  instances over 24 hours, with one planned maintenance event (SURG-01
  maintenance window from 22:15 to 01:59) managed without patient
  impact~\cite{newtrial2026}.

\item Sim-to-real divergence exceeding 4 millimeters in trajectory or 1.0
  Newton in force shall trigger an automatic safety review and potential
  suspension of the affected system~\cite{cfr312adapt}.
\end{enumerate}

\subsection{Adverse Event Response}

\begin{enumerate}[label=(\alph*)]
\item Every Physical AI-related adverse event shall be documented, reported, and
  resolved in accordance with the timelines and categories established in the
  adapted 21 CFR Part 312 \S~312.32~\cite{cfr312adapt}.

\item The site shall achieve a target of resolving all Grade 1 and Grade 2
  adverse events within the same hour of occurrence, consistent with the
  performance demonstrated in the simulation where all seven adverse events
  (ranging from nausea and hypotension to pain, cough, desaturation, and
  anxiety) were resolved within the hour~\cite{newtrial2026}.

\item A Physical AI Safety Review Committee shall convene within 24 hours of any
  Grade 3 or higher adverse event and within 7 days of any fatal or
  life-threatening event.

\item The committee shall include expertise in robotic engineering, AI and
  machine learning, oncology, and cybersecurity, meeting at regular intervals
  not exceeding 90 days~\cite{cfr312adapt}.
\end{enumerate}

\section{Patient Complaint and Grievance Process}

\begin{enumerate}[label=(\alph*)]
\item Each authorized site shall establish a patient complaint and grievance
  process that is accessible 24 hours per day, 7 days per week, consistent
  with the site's operating schedule.

\item Complaints involving Physical AI system malfunctions or safety concerns
  shall be investigated within 48 hours and the patient shall be notified of
  the investigation outcome within 14 days.

\item The department shall maintain a public registry of Physical AI-related
  patient complaints and their resolutions, with appropriate de-identification
  of patient information.

\item Patients shall have the right to file complaints directly with the
  department if they believe the site's internal grievance process has not
  adequately addressed their concerns.
\end{enumerate}

\section{Enforcement}

\begin{enumerate}[label=(\alph*)]
\item The department shall have authority to inspect authorized Physical AI
  oncology clinical trial sites without prior notice to verify compliance
  with this article.

\item Violations of patient rights established under this article shall be
  subject to civil penalties of not less than one thousand dollars (\$1,000)
  and not more than twenty-five thousand dollars (\$25,000) per violation.

\item The department may suspend or revoke site authorization for repeated or
  serious violations of patient rights or robotic safety standards.

\item Nothing in this article shall be construed to limit any other remedy
  available to patients under California law.
\end{enumerate}

\section{Effective Date}

This act shall take effect on January 1, 2027.

\printbibliography[title={References}]

\end{document}
